\documentclass[aspectratio=169]{ctexbeamer}

\usepackage{amsmath}
\usepackage{booktabs}
\usepackage{graphicx}
\usepackage{listings}
\usepackage{tikz}
\usepackage{xcolor}
\usetikzlibrary{arrows.meta, positioning, shapes.geometric}

\usetheme{Berlin}
\usecolortheme{default}

\definecolor{mainblue}{RGB}{32, 83, 154}
\definecolor{deepblue}{RGB}{18, 55, 116}
\definecolor{footblue}{RGB}{38, 92, 166}
\definecolor{softblue}{RGB}{230, 238, 250}
\definecolor{softgreen}{RGB}{232, 245, 233}
\definecolor{softorange}{RGB}{255, 243, 224}
\definecolor{softred}{RGB}{252, 232, 232}
\definecolor{codebg}{RGB}{248, 248, 246}
\definecolor{codekw}{RGB}{126, 63, 160}
\definecolor{codecomment}{RGB}{70, 130, 70}
\definecolor{codered}{RGB}{160, 60, 60}

\setbeamercolor{structure}{fg=mainblue}
\setbeamercolor{block title}{fg=white,bg=mainblue}
\setbeamercolor{block body}{bg=softblue}
\setbeamercolor{foot left}{fg=white,bg=footblue}
\setbeamercolor{foot middle}{fg=white,bg=deepblue}
\setbeamercolor{foot right}{fg=white,bg=footblue}
\setbeamertemplate{navigation symbols}{}

% Berlin/miniframes 默认包含 section 与 subsection 两条顶部导航；
% 本报告没有 subsection，覆盖模板以去掉空白导航行。
\setbeamertemplate{headline}{%
  \begin{beamercolorbox}[colsep=1.2pt]{upper separation line head}
  \end{beamercolorbox}
  \begin{beamercolorbox}{section in head/foot}
    \vskip2.3ex
    \vskip-\baselineskip
    \insertnavigation{\paperwidth}%
    \vskip1.0ex
  \end{beamercolorbox}%
  \begin{beamercolorbox}[colsep=1.2pt]{lower separation line head}
  \end{beamercolorbox}
}

% 将 Berlin 默认的两行页脚合并为一行，并用三栏颜色区分。
\setbeamertemplate{footline}{%
  \leavevmode%
  \begin{beamercolorbox}[wd=.33\paperwidth,ht=2.8ex,dp=1ex,leftskip=.32cm]{foot left}
    \usebeamerfont{author in head/foot}\insertshortauthor
  \end{beamercolorbox}%
  \begin{beamercolorbox}[wd=.34\paperwidth,ht=2.8ex,dp=1ex,center]{foot middle}
    \usebeamerfont{title in head/foot}编译原理大作业
  \end{beamercolorbox}%
  \begin{beamercolorbox}[wd=.33\paperwidth,ht=2.8ex,dp=1ex,rightskip=.32cm plus1fil]{foot right}
    \hfill\usebeamerfont{page number in head/foot}\insertframenumber{} / \inserttotalframenumber
  \end{beamercolorbox}%
}

\lstdefinelanguage{RustLike}{
    morekeywords={fn, let, mut, if, else, while, for, in, loop, break, continue, return, i32},
    sensitive=true,
    morecomment=[l]{//},
    morecomment=[s]{/*}{*/},
    morestring=[b]"
}

\lstdefinestyle{code}{
    language=RustLike,
    backgroundcolor=\color{codebg},
    basicstyle=\ttfamily\scriptsize,
    keywordstyle=\color{codekw}\bfseries,
    commentstyle=\color{codecomment},
    stringstyle=\color{codered},
    showstringspaces=false,
    breaklines=true,
    columns=fullflexible,
    keepspaces=true,
    frame=single,
    rulecolor=\color{black!18},
    xleftmargin=0.45em,
    xrightmargin=0.45em
}

\lstdefinestyle{tinycode}{
    style=code,
    basicstyle=\ttfamily\tiny
}

\title[编译原理大作业 1 \& 2]{编译原理大作业 1 \& 2}
\subtitle{类 Rust 语言的词法/语法分析、语义分析与四元式中间代码生成}
\author[李堂辉, 高胜寒, 冉子易]{2351078 李堂辉 \\ 2351837 高胜寒 \\ 2352197 冉子易\\
\small （按学号为序）}
\institute{同济大学}
\date{2026年6月16日}

\begin{document}

\begin{frame}
    \titlepage
\end{frame}

\begin{frame}{主线：两个作业串成一条前端流程}
    \begin{center}
    \resizebox{0.96\linewidth}{!}{
    \begin{tikzpicture}[
        node distance=1.25cm,
        box/.style={rectangle, rounded corners, draw=mainblue, thick, minimum width=2.55cm, minimum height=0.78cm, align=center, fill=softblue},
        hot/.style={rectangle, rounded corners, draw=mainblue, thick, minimum width=2.55cm, minimum height=0.78cm, align=center, fill=softorange},
        outbox/.style={rectangle, rounded corners, draw=mainblue, thick, minimum width=2.55cm, minimum height=0.78cm, align=center, fill=softgreen},
        arr/.style={-{Latex[length=2.8mm]}, thick, draw=mainblue}
    ]
        \node[box] (src) {类 Rust\\源程序};
        \node[box, right=of src] (tok) {Token\\带行列号};
        \node[hot, right=of tok] (ast) {AST\\结构化契约};
        \node[box, right=of ast] (sem) {语义状态\\作用域/类型/循环};
        \node[outbox, right=of sem] (ir) {四元式 IR\\线性代码};
        \draw[arr] (src) -- (tok);
        \draw[arr] (tok) -- (ast);
        \draw[arr] (ast) -- (sem);
        \draw[arr] (sem) -- (ir);
        \node[below=0.28cm of tok, text=mainblue] {\small 大作业 1};
        \node[below=0.28cm of sem, text=mainblue] {\small 大作业 2};
    \end{tikzpicture}}
    \end{center}
    \vspace{-0.1em}
    \begin{block}{实现重点}
        这份实现把类 Rust 语法保留为结构化 AST，并在第二个作业中继续完成语义诊断和 IR 生成。
    \end{block}
\end{frame}

\begin{frame}{对标要求：覆盖情况与实现亮点}
    \begin{table}
    \centering
    \small
    \begin{tabular}{p{2.9cm}p{4.4cm}p{4.4cm}}
        \toprule
        \textbf{要求项} & \textbf{完成情况} & \textbf{实现亮点} \\
        \midrule
        大作业 1 必做语法 & 规则覆盖到 5.1 & 递归下降解析和 AST 构造 \\
        大作业 1 扩展结构 & 引用、数组、元组、表达式块、三类循环 & 统一保留到 AST 节点 \\
        大作业 2 语义诊断 & 类型、可变性、返回、循环控制检查 & 符号表栈、循环栈、错误汇总 \\
        大作业 2 中间代码 & 生成四元式 IR & 临时变量、标签和跳转 \\
        \bottomrule
    \end{tabular}
    \end{table}
    \vspace{0.2em}
    \centering
    两个作业共用 AST 作为衔接点，语义检查和 IR 生成在其上继续展开。
\end{frame}

\begin{frame}{实现特色：AST 作为两份作业的接口}
    \begin{table}
    \centering
    \small
    \begin{tabular}{p{3.2cm}p{4.1cm}p{4.6cm}}
        \toprule
        \textbf{作业 1 保留的信息} & \textbf{代码落点} & \textbf{作业 2 的使用方式} \\
        \midrule
        类型结构 &
        \begin{tabular}[t]{@{}l@{}}\texttt{Ref/MutRef}\\\texttt{Array/Tuple}\end{tabular} &
        初始化、赋值、返回值和索引检查 \\
        左值形态 &
        \begin{tabular}[t]{@{}l@{}}\texttt{Ident/Deref}\\\texttt{Index/TupleIndex}\end{tabular} &
        区分变量赋值、解引用赋值和聚合访问 \\
        控制流结构 &
        \begin{tabular}[t]{@{}l@{}}\texttt{If/While/For/Loop}\\\texttt{Break/Continue}\end{tabular} &
        生成标签，并维护 \texttt{loop\_stack} \\
        表达式结构 &
        \begin{tabular}[t]{@{}l@{}}\texttt{ExprBlock/IfExpr}\\\texttt{Call}\end{tabular} &
        返回 \texttt{(Type, Arg)} 继续生成 IR \\
        \bottomrule
    \end{tabular}
    \end{table}
    \vspace{0.2em}
    \begin{block}{关键点}
        AST 在两个工程之间承担稳定接口，语义分析和 IR 生成只处理结构化节点。
    \end{block}
\end{frame}

\section{大作业 1：把类 Rust 代码变成结构}

\begin{frame}{作业 1 特色：扩展语法保留成可分析结构}
    \begin{columns}[T]
    \begin{column}{0.50\textwidth}
        \begin{block}{覆盖范围}
            \scriptsize
            \begin{tabular}{p{1.2cm}p{4.2cm}}
                类型 & \texttt{i32}, \texttt{\&T}, \texttt{\&mut T}, 数组, 元组 \\
                左值 & 标识符, 解引用, 数组下标, 元组下标 \\
                控制流 & \texttt{if}, \texttt{while}, \texttt{for}, \texttt{loop} \\
                表达式 & 调用, 区间, 表达式块, if 表达式 \\
            \end{tabular}
        \end{block}
    \end{column}
    \begin{column}{0.46\textwidth}
        \begin{itemize}
            \item 引用、可变引用、数组、元组直接进入 \texttt{Type} 和 \texttt{Expr}。
            \item 表达式块、\texttt{if} 表达式保留“表达式有值”的语言特征。
            \item \texttt{while/for/loop} 与 \texttt{break/continue} 为循环栈提供结构入口。
        \end{itemize}
    \end{column}
    \end{columns}
\end{frame}

\begin{frame}{Parser：把复杂语法压进可维护的层次}
    \begin{center}
    \resizebox{0.98\linewidth}{!}{
    \begin{tikzpicture}[
        node distance=0.72cm,
        box/.style={rectangle, rounded corners, draw=mainblue, thick, minimum width=3.0cm, minimum height=0.62cm, align=center, fill=softblue},
        hot/.style={rectangle, rounded corners, draw=mainblue, thick, minimum width=3.0cm, minimum height=0.62cm, align=center, fill=softorange},
        arr/.style={-{Latex[length=2.4mm]}, thick, draw=mainblue}
    ]
        \node[box] (program) {\texttt{parse\_program}};
        \node[box, below=of program] (decl) {\texttt{parse\_function\_decl}};
        \node[box, below left=0.62cm and 1.8cm of decl] (block) {\texttt{parse\_block}};
        \node[box, below right=0.62cm and 1.8cm of decl] (type) {\texttt{parse\_type}};
        \node[box, below=of block] (stmt) {\texttt{parse\_statement}};
        \node[hot, below=of stmt] (expr) {\texttt{parse\_expression}};
        \node[hot, below left=0.62cm and 1.0cm of expr] (bin) {区间/比较/加减/乘除};
        \node[hot, below right=0.62cm and 1.0cm of expr] (post) {调用/索引/字段/引用};
        \draw[arr] (program) -- (decl);
        \draw[arr] (decl) -- (block);
        \draw[arr] (decl) -- (type);
        \draw[arr] (block) -- (stmt);
        \draw[arr] (stmt) -- (expr);
        \draw[arr] (expr) -- (bin);
        \draw[arr] (expr) -- (post);
    \end{tikzpicture}}
    \end{center}
    \scriptsize
    这里保留递归下降的优势：结构规则和表达式优先级分开，出错位置也容易落到具体解析函数。
\end{frame}

\begin{frame}[fragile]{AST：保留语义和 IR 需要的信息}
    \begin{columns}[T]
    \begin{column}{0.48\textwidth}
        \begin{lstlisting}[style=tinycode]
pub enum Type {
    I32,
    Ref(Box<Type>),
    MutRef(Box<Type>),
    Array(Box<Type>, i64),
    Tuple(Vec<Type>),
}
        \end{lstlisting}
    \end{column}
    \begin{column}{0.48\textwidth}
        \begin{lstlisting}[style=tinycode]
pub enum Expr {
    BinaryOp(...),
    Ref(Box<Expr>),
    MutRef(Box<Expr>),
    Index(Box<Expr>, Box<Expr>),
    TupleIndex(Box<Expr>, i64),
    ExprBlock(ExprBlock),
    IfExpr(...),
}
        \end{lstlisting}
    \end{column}
    \end{columns}
    \vspace{0.2em}
    \begin{block}{设计取向}
        AST 保留类型、左值形态、控制流和表达式形态，给第二个作业直接消费。
    \end{block}
\end{frame}

\begin{frame}{作业 1 产出：可继续编译的结构}
    \begin{table}
    \centering
    \small
    \begin{tabular}{p{2.5cm}p{5.1cm}p{4.9cm}}
        \toprule
        \textbf{产物} & \textbf{包含的信息} & \textbf{给作业 2 的价值} \\
        \midrule
        Token & 类别、字面量、行列号 & 定位词法/语法问题 \\
        AST 语句 & 声明、赋值、返回、分支、循环 & 语义检查按语句类型分派 \\
        AST 表达式 & 类型相关结构、调用、索引、引用 & 推导类型并生成临时变量 \\
        AST 类型 & 五类 \texttt{Type} & 赋值、返回、解引用检查的依据 \\
        \bottomrule
    \end{tabular}
    \end{table}
    \vspace{0.2em}
    \centering
    这就是两个大作业能串起来的关键。
\end{frame}

\section{大作业 2：让 AST 带上语义状态}

\begin{frame}{作业 2 特色：带状态遍历 AST}
    \begin{center}
    \resizebox{0.92\linewidth}{!}{
    \begin{tikzpicture}[
        node distance=1.2cm,
        box/.style={rectangle, rounded corners, draw=mainblue, thick, minimum width=2.65cm, minimum height=0.78cm, align=center, fill=softblue},
        state/.style={rectangle, rounded corners, draw=mainblue, thick, minimum width=2.65cm, minimum height=0.78cm, align=center, fill=softorange},
        outbox/.style={rectangle, rounded corners, draw=mainblue, thick, minimum width=2.65cm, minimum height=0.78cm, align=center, fill=softgreen},
        arr/.style={-{Latex[length=2.8mm]}, thick, draw=mainblue}
    ]
        \node[box] (ast) {AST};
        \node[state, right=of ast] (state) {\texttt{Compiler}\\状态机};
        \node[outbox, right=of state] (ok) {合法：\\IR};
        \node[outbox, below=0.62cm of ok] (err) {非法：\\错误列表};
        \draw[arr] (ast) -- (state);
        \draw[arr] (state) -- (ok);
        \draw[arr] (state) -- (err);
    \end{tikzpicture}}
    \end{center}
    \begin{table}
    \centering
    \scriptsize
    \begin{tabular}{p{2.7cm}p{3.7cm}p{4.7cm}}
        \toprule
        \textbf{状态} & \textbf{代码字段} & \textbf{用途} \\
        \midrule
        作用域 & \texttt{scopes} & 找变量、查类型、查可变性 \\
        函数返回 & \texttt{current\_return\_type} & 检查 \texttt{return} 是否匹配签名 \\
        循环上下文 & \texttt{loop\_stack} & 检查并翻译 \texttt{break/continue} \\
        生成计数 & \texttt{temp\_count, label\_count} & 分配 \texttt{tN} 和 \texttt{LN} \\
        \bottomrule
    \end{tabular}
    \end{table}
\end{frame}

\begin{frame}[fragile]{符号表栈：作用域、类型、可变性放在一起}
    \begin{columns}[T]
    \begin{column}{0.49\textwidth}
        \begin{lstlisting}[style=tinycode]
struct SymbolInfo {
    mutable: bool,
    ty: Type,
    arg: Arg,
}

struct Compiler {
    scopes: Vec<HashMap<String, SymbolInfo>>,
    errors: Vec<String>,
}
        \end{lstlisting}
    \end{column}
    \begin{column}{0.47\textwidth}
        \begin{itemize}
            \item 进入函数、块、表达式块时压栈。
            \item 退出作用域时弹栈，局部变量自然失效。
            \item 查找变量时从内层向外层倒序查找。
            \item 允许 shadowing，和类 Rust 语言直觉一致。
        \end{itemize}
    \end{column}
    \end{columns}
\end{frame}

\begin{frame}{语义检查：回应要求 PPTX 的易错点}
    \begin{table}
    \centering
    \small
    \begin{tabular}{p{2.9cm}p{4.1cm}p{4.4cm}}
        \toprule
        \textbf{检查点} & \textbf{触发场景} & \textbf{背后的状态} \\
        \midrule
        未声明变量 & \texttt{a = 1;} 但没有声明 & \texttt{lookup\_variable} \\
        不可变变量赋值 & \texttt{let a = 1; a = 2;} & \texttt{mutable} 字段 \\
        类型不匹配 & 数组赋给 \texttt{i32}，或 \texttt{i32 + tuple} & \texttt{Type} 枚举 \\
        引用/解引用约束 & 读解引用、写解引用、非引用解引用 & \texttt{Ref/MutRef} \\
        返回类型 & \texttt{-> i32} 但 \texttt{return;} & 当前函数返回类型 \\
        循环控制 & 循环外 \texttt{break/continue} & \texttt{loop\_stack} \\
        \bottomrule
    \end{tabular}
    \end{table}
\end{frame}

\begin{frame}[fragile]{错误汇总：一次运行看出多类问题}
    \begin{columns}[T]
    \begin{column}{0.45\textwidth}
        \begin{lstlisting}[style=tinycode]
fn test() -> i32 {
    let a: i32 = 1;
    a = 2;
    break;
    return;
}
        \end{lstlisting}
    \end{column}
    \begin{column}{0.51\textwidth}
        \begin{lstlisting}[style=tinycode]
[Semantic Error] 不可变变量 'a' 不允许被二次赋值
[Semantic Error] break 必须出现在循环体内
[Semantic Error] 返回语句类型空和函数声明类型 I32 不一致
        \end{lstlisting}
    \end{column}
    \end{columns}
    \vspace{0.2em}
    \begin{block}{实现方式}
        \texttt{Compiler} 会继续收集语义错误，统一写入 \texttt{errors: Vec<String>}，最后集中输出。
    \end{block}
\end{frame}

\section{大作业 2：把结构化程序压平成 IR}

\begin{frame}[fragile]{IR：从树形结构变成线性四元式}
    \begin{columns}[T]
    \begin{column}{0.48\textwidth}
        \begin{lstlisting}[style=tinycode]
pub enum Arg {
    Empty,
    Int(i64),
    Var(String),
    Temp(usize),
    Label(usize),
}
        \end{lstlisting}
    \end{column}
    \begin{column}{0.48\textwidth}
        \begin{lstlisting}[style=tinycode]
pub struct Quadruple {
    op: Op,
    arg1: Arg,
    arg2: Arg,
    result: Arg,
}
        \end{lstlisting}
    \end{column}
    \end{columns}
    \begin{table}
    \centering
    \scriptsize
    \begin{tabular}{p{2.5cm}p{7.8cm}}
        \toprule
        \textbf{指令族} & \textbf{代表指令} \\
        \midrule
        算术/比较 & \texttt{ADD, SUB, MUL, DIV, EQ, NE, LT, LE, GT, GE} \\
        控制流 & \texttt{LABEL, JUMP, JUMP\_IF\_FALSE} \\
        调用/返回 & \texttt{ARG, CALL, PARAM, RETURN} \\
        引用/聚合 & \texttt{REF, DEREF, DEREF\_ASSIGN, LOAD\_ARRAY} \\
        \bottomrule
    \end{tabular}
    \end{table}
\end{frame}

\begin{frame}[fragile]{表达式翻译：临时变量串起计算}
    \begin{columns}[T]
    \begin{column}{0.39\textwidth}
        \begin{lstlisting}[style=code]
let mut c = a + b * 2;
        \end{lstlisting}
        \vspace{0.4em}
        \small
        递归翻译表达式时，每个中间结果都拿到一个 \texttt{tN}。
    \end{column}
    \begin{column}{0.57\textwidth}
        \begin{lstlisting}[style=code]
(MUL,    b,   2,   t0)
(ADD,    a,  t0,   t1)
(ASSIGN, t1,  -,    c)
        \end{lstlisting}
    \end{column}
    \end{columns}
    \vspace{0.3em}
    \centering
    \texttt{compile\_expression} 的返回值是 \texttt{(Type, Arg)}：既给语义检查，也给 IR 继续使用。
\end{frame}

\begin{frame}[fragile]{分支翻译：条件失败跳转}
    \begin{columns}[T]
    \begin{column}{0.39\textwidth}
        \begin{lstlisting}[style=tinycode]
if a > 0 {
    return 1;
} else {
    return 0;
}
        \end{lstlisting}
    \end{column}
    \begin{column}{0.57\textwidth}
        \begin{lstlisting}[style=tinycode]
(GT,             a, 0, t0)
(JUMP_IF_FALSE, t0, -, L0)
(RETURN,         1, -, -)
(JUMP,           -, -, L1)
L0:
(RETURN,         0, -, -)
L1:
        \end{lstlisting}
    \end{column}
    \end{columns}
    \vspace{0.2em}
    \centering
    结构化的 \texttt{if/else} 被压成“条件、失败标签、结束标签”。
\end{frame}

\begin{frame}{循环翻译：标签和循环栈}
    \begin{columns}[T]
    \begin{column}{0.45\textwidth}
        \begin{block}{入口标签}
            \centering
            \texttt{start\_label}
        \end{block}
        \begin{itemize}
            \item \texttt{while} 重新判断条件。
            \item \texttt{loop} 直接进入下一轮。
            \item \texttt{continue} 跳到这里。
        \end{itemize}
    \end{column}
    \begin{column}{0.45\textwidth}
        \begin{block}{出口标签}
            \centering
            \texttt{end\_label}
        \end{block}
        \begin{itemize}
            \item 条件失败跳到这里。
            \item \texttt{break} 跳到这里。
            \item 循环体结束后接下一条指令。
        \end{itemize}
    \end{column}
    \end{columns}
    \vspace{0.3em}
    \begin{center}
        \texttt{loop\_stack.push((start\_label, end\_label, None))}
    \end{center}
\end{frame}

\begin{frame}{扩展结构：语义和 IR 都有落点}
    \begin{table}
    \centering
    \small
    \begin{tabular}{p{2.8cm}p{4.0cm}p{4.6cm}}
        \toprule
        \textbf{源码结构} & \textbf{AST 节点} & \textbf{IR 或语义落点} \\
        \midrule
        \texttt{\&x / \&mut x} & \texttt{Ref / MutRef} & 生成 \texttt{REF}，类型变成引用类型 \\
        \texttt{*p} & \texttt{UnaryDeref} & 检查引用类型，生成 \texttt{DEREF} \\
        \texttt{a[i]} & \texttt{Index} & 检查数组类型，读取生成 \texttt{LOAD\_ARRAY} \\
        \texttt{t.0} & \texttt{TupleIndex} & 根据元组类型推导元素类型 \\
        \texttt{foo(a,b)} & \texttt{Call} & 参数生成 \texttt{ARG}，调用生成 \texttt{CALL} \\
        \bottomrule
    \end{tabular}
    \end{table}
    \vspace{0.2em}
    \small
    当前实现把这些结构接到语义/IR 主流程；函数签名检查和数组写入地址计算可作为后续扩展。
\end{frame}

\section{验证与收束}

\begin{frame}{测试覆盖：功能集和诊断集分开}
    \begin{table}
    \centering
    \scriptsize
    \begin{tabular}{p{2.7cm}p{4.1cm}p{4.7cm}}
        \toprule
        \textbf{测试入口} & \textbf{覆盖内容} & \textbf{观察结果} \\
        \midrule
        \texttt{test\_all.rs.txt} & 函数、声明、表达式、分支、循环、引用、数组、元组 & 作业 1 输出 Token/AST，作业 2 输出 IR \\
        \texttt{test\_errors.rs.txt} & 非法字符、缺分号、括号不匹配等 & Lexer/Parser 报错 \\
        \texttt{test\_semantic}\\\texttt{\_errors.rs.txt} & 未声明、类型不匹配、不可变赋值、非法控制流、返回类型错误 & \texttt{[Semantic Error]} 列表 \\
        \texttt{cargo check} & Rust 工程可编译性 & 两个正式工程均无编译错误 \\
        \bottomrule
    \end{tabular}
    \end{table}
\end{frame}

\begin{frame}{汇报收束：我们真正完成的是一条前端流水线}
    \begin{block}{一句话总结}
        这两个作业合起来完成了从类 Rust 源程序到四元式 IR 的前端流程。
    \end{block}
    \begin{itemize}
        \item 作业 1 的重点是把复杂语法结构化，尤其是引用、数组、元组、表达式块和循环结构。
        \item 作业 2 的重点是带状态遍历 AST，把作用域、类型、可变性、返回类型和循环上下文检查清楚。
        \item 最终 IR 用临时变量和标签把表达式和控制流线性化，为后续解释执行或目标代码生成留下接口。
    \end{itemize}
    \vfill
    \centering
    \Large 谢谢老师，欢迎提问
\end{frame}

\appendix
\section{备用}

\begin{frame}{备用：如果问“两份作业区别是什么？”}
    \begin{table}
    \centering
    \small
    \begin{tabular}{p{2.3cm}p{3.0cm}p{3.2cm}p{3.1cm}}
        \toprule
        \textbf{项目} & \textbf{输入} & \textbf{核心状态} & \textbf{输出} \\
        \midrule
        大作业 1 & 源代码字符串 & Token 指针、解析函数栈 & Token + AST \\
        大作业 2 & AST & 符号表栈、循环栈、返回类型、计数器 & 语义错误或 IR \\
        \bottomrule
    \end{tabular}
    \end{table}
\end{frame}

\begin{frame}[fragile]{备用：while 的 IR 生成细节}
    \begin{lstlisting}[style=tinycode]
let start_label = self.new_label();
let end_label = self.new_label();

emit(Label,       -,        -, Label(start_label));
let cond_arg = compile_expression(cond);
emit(JumpIfFalse, cond_arg, -, Label(end_label));

loop_stack.push((start_label, end_label, None));
compile_block(body);
loop_stack.pop();

emit(Jump,        -,        -, Label(start_label));
emit(Label,       -,        -, Label(end_label));
    \end{lstlisting}
\end{frame}

\begin{frame}{备用：实现边界}
    \begin{itemize}
        \item 当前重点是前端流程和四元式表达，没有继续做目标代码生成。
        \item 函数调用已经生成 \texttt{ARG/CALL}，但函数签名检查仍可继续完善。
        \item 数组读取有 \texttt{LOAD\_ARRAY}，数组写入接口存在，后续可以扩展更完整的地址计算和 \texttt{STORE\_ARRAY}。
        \item \texttt{for/range} 进入语法树，IR 层当前重点展示 \texttt{while/loop/break/continue} 的标签化控制流。
    \end{itemize}
\end{frame}

\end{document}
